\section{Holomorphic operator substitution}\label{sec:operator_substitution}

The scalar identities have a direct functional calculus form.  This section records only the basic consequence and does not develop an operator resolvent theory.

\begin{standardfact}[Holomorphic functional calculus]\label{sf:holomorphic_calculus}
Let $A$ be a bounded operator on a complex Banach space.  If $f$ is holomorphic on a neighborhood of $\Spec(A)$, then
\[
f(A)
=
\frac1{2\pi i}\int_\Gamma f(z)(zI-A)^{-1}\,dz
\]
for a contour $\Gamma$ surrounding the spectrum.  The calculus preserves sums, products, and composition.  We use the standard form from \cite{DunfordSchwartz1958}.
\end{standardfact}

\begin{proposition}[Operator reciprocal family]\label{prop:operator_reciprocal}
Let $A$ be bounded and suppose $\Spec(A)$ lies in a simply connected region on which $\Dfun_N$ has no zeros.  Fix a holomorphic branch $g=\log\Dfun_N$ there.  Then
\[
\Sfun_N^{(r)}(A)
=
\exp\bigl(-r g(A)\bigr)
=
\bigl(\Dfun_N^{-r}\bigr)(A)
\]
is well defined and
\[
\Sfun_N^{(r)}(A)\Sfun_N^{(s)}(A)
=
\Sfun_N^{(r+s)}(A).
\]
If $(P_i)$ is a finite complete family of bounded idempotents, then
\[
P_i f(A)P_j
=
\frac1{2\pi i}
\int_\Gamma
f(z)P_i(zI-A)^{-1}P_j\,dz
\]
for every holomorphic $f$ on a neighborhood of the spectrum.
\end{proposition}

\begin{proof}
The spectral mapping theorem makes $\Dfun_N(A)$ invertible.  The product identity follows from multiplicativity of the functional calculus and the fixed logarithm.  Applying $P_i$ and $P_j$ to the Cauchy integral formula gives the packet formula.
\end{proof}

